\section{Related Work}
\label{sec:related-work}

Research related to Stockrium lies at the intersection of financial decision support, large language models, multi-agent systems, financial-news sentiment analysis, and Vietnamese stock-market prediction. This section reviews representative studies in these areas and identifies the gap addressed by the proposed application.

\subsection{Multi-Agent LLM Systems for Financial Decision Support}

Multi-agent LLM systems divide a complex financial task among components with specialized roles, allowing technical, fundamental, news, trading, and risk perspectives to be processed separately. TradingAgents is a representative framework that models a trading firm through analyst agents, Bull and Bear researchers, a trader, and a risk-management team \cite{xiao2024tradingagents}. Its reported results demonstrate that structured collaboration can combine heterogeneous evidence and improve selected return and risk measures. However, its multi-round debates require numerous model calls, increase latency and cost, and leave the final action strongly dependent on the language model's reasoning path. FinMem instead equips a single trading agent with layered memory that retrieves financial information according to temporal relevance \cite{yu2023finmem}, highlighting the value of maintaining context across changing market conditions, although its main objective remains autonomous trading performance. Stockrium adopts role specialization but uses a smaller graph designed for Vietnamese equities: technical, fundamental, and news agents operate in parallel, after which their evidence is synthesized. Deterministic scoring and validation are added around the LLM output to improve traceability and constrain the final recommendation. Thus, Stockrium is positioned as a user-facing decision-support system rather than an autonomous trading agent or a debate-based simulation of a financial institution.

\subsection{Finance-Specific Large Language Models}

Finance-specific language models show that domain knowledge can improve the interpretation of specialized terminology, disclosures, and market communication. BloombergGPT is a 50-billion-parameter model trained on both financial and general-purpose corpora; its evaluation indicates that domain-focused pretraining improves performance on financial NLP tasks while retaining general capabilities \cite{wu2023bloomberggpt}. FinGPT provides a more accessible, open-source alternative centered on automated financial-data curation and lightweight model adaptation for tasks such as sentiment analysis, robo-advising, and algorithmic trading \cite{yang2023fingpt}. These studies establish the usefulness of LLMs for extracting information, understanding financial text, and generating explanations. Nevertheless, a financial language model alone does not constitute a complete decision-support application: it does not necessarily retrieve synchronized market and company data, reconcile multiple analytical perspectives, apply user configurations, enforce numerical risk constraints, or evaluate strategies historically. The cited models are also not specifically designed for Vietnamese financial language or local market conventions. Consequently, Stockrium uses general LLM capabilities within a controlled application workflow, while structured data processing and deterministic rules remain responsible for calculations and output validation.

\subsection{Financial News Sentiment Analysis}

Financial news may contain events and expectations that are not fully reflected in historical prices, but its specialized vocabulary limits the effectiveness of general sentiment models. FinBERT addresses this problem through domain-specific pretraining and fine-tuning, achieving improved positive, neutral, and negative classification on financial datasets \cite{araci2019finbert}. FinBERT-LSTM subsequently combined news sentiment with historical prices for NASDAQ-100 prediction, showing that textual information can supplement numerical time series \cite{halder2022finbertlstm}; however, it remains a price-prediction pipeline and does not reconcile news with fundamental and technical evidence. Vietnamese studies provide more locally relevant findings. Tung et al. reported that PhoBERT outperformed conventional models on Vietnamese stock-article titles, reaching approximately 93\% accuracy \cite{tung2021phobert}. Using a much larger collection of financial and economic articles, Vu et al. achieved over 81\% sentiment-classification accuracy and found that negative news was associated more clearly with return variance than with a stable change in mean return \cite{vu2023sentiments}. These results suggest that sentiment can provide useful contextual and risk information but should not independently determine an investment action. Stockrium therefore summarizes and assesses stock-related news as one analytical source whose output is combined with technical and fundamental evidence.

\subsection{Prediction Research for the Vietnamese Stock Market}

Vietnamese stock-market research has primarily used historical numerical data to forecast index or stock movements. Truong et al. compared LSTM variants for VN-Index prediction and found bidirectional LSTM effective across the examined horizons \cite{truong2022lstm}. Phuoc et al. incorporated moving averages, MACD, and RSI into an LSTM-based model for the VN-Index and VN30 stocks, reporting high predictive accuracy for most analyzed series \cite{phuoc2024ml}. Son et al. later compared multilayer perceptron, LSTM, bidirectional LSTM, and Transformer models and obtained horizon-dependent results: Transformer performed better for short-term forecasting, whereas bidirectional LSTM was stronger over medium and longer horizons \cite{son2025deepmodels}. These studies confirm the relevance of machine learning, technical indicators, and investment horizon in the Vietnamese market. Nevertheless, their main evaluation target is forecast error or directional accuracy, which does not directly indicate whether a usable strategy controls drawdown, remains effective after transaction costs, or provides understandable evidence for a recommendation. They also rarely combine price behavior with company fundamentals and Vietnamese news. Stockrium extends this research direction by emphasizing integrated and explainable decision support together with historical strategy evaluation rather than price prediction alone.

\subsection{Research Gap and Positioning of Stockrium}

The literature reveals a separation among financial AI research streams: forecasting studies focus mainly on prices and technical indicators, sentiment studies classify news, finance-specific LLMs address language understanding, and multi-agent frameworks often target autonomous trading in international markets. Few systems combine technical, fundamental, and Vietnamese news evidence in a user-facing application while also supporting configurable analysis, investment-horizon awareness, explainable recommendations, deterministic decision constraints, and backtesting. Stockrium is designed to address this integration gap. It coordinates specialized agents through an explicit workflow, allows users to select analytical inputs and their relative importance, and uses an LLM to interpret and synthesize evidence rather than perform every calculation. Deterministic components calculate indicators and confidence, validate minimum conditions, and restrict unsupported recommendations, while the backtesting facility enables the resulting strategy to be examined on historical Vietnamese data. The contribution is therefore not a new forecasting model or a claim of guaranteed superior returns; it is the construction of a localized, configurable, and transparent decision-support environment that brings previously separate analytical capabilities into a coherent system.
